% BEGIN_REPORT::LEAD::ARC-0007::2026-10-07T23:34:47-07:00
\documentclass[11pt]{article}
\usepackage[T1]{fontenc}
\usepackage[margin=.86in]{geometry}
\usepackage{amsmath,amssymb,amsthm,booktabs,longtable,xurl,hyperref}
\setlength{\emergencystretch}{3em}
\newtheorem{theorem}{Theorem}
\newtheorem{definition}[theorem]{Definition}
\newcommand{\SAT}{\operatorname{SAT}}
\newcommand{\R}{\mathcal R}
\newcommand{\B}{\mathcal B}
\title{ARC-0007: Finite-State Phase Potentials\\for Existential Achievable-Region Closure}
\author{LEAD, ARC/WRAPPER research team}
\date{8 October 2026}
\begin{document}
\maketitle

\begin{center}
\fbox{\parbox{.92\linewidth}{
\textbf{Result.}
PHASE-MULTINORM-ARC1 is \textsc{Proved/Exact} as a conditional
achievable-region envelope theorem.  It permits a fixed finite graph of
phase-dependent positive linear potentials and existentially chosen successor
edges.  A two-phase exact rational certificate works although its two cross
matrices are individually unstable and no common positive linear potential
can contract them.  The bounded regression is \textsc{Computational}.
No SAT-specific edge relation was constructed.  That implication, the
model-matched endpoint, and P versus NP remain \textsc{Open}.}}
\end{center}

\section{Claim card}
\begin{description}
\item[] \textbf{Claim ID:} PHASE-MULTINORM-ARC1.
\item[] \textbf{Status before run:} fixed affine matrices, prescribed finite
return words, and cofinal finite routes were already conditional
\textsc{Proved/Exact}; a phase-dependent graph potential was absent.
\item[] \textbf{Desired conclusion:} a finite achieved exponent pair enters
\(\B=\{(t,s)\ge0:t+s<\gamma\}\), where
\(\gamma=2\cos(\pi/7)\), from graph-local affine successor bounds.
\item[] \textbf{Exact quantifier prefix:}
\[
\exists\text{ fixed finite }G,\{M_e,d_e\},\{v_q\},\lambda,c,\theta\
\forall q\in Q\ \forall z\in\R\setminus\B\
\exists e:q\to r\ \exists z'\in\R:
z'\le M_ez+d_e.
\]
The edge inequalities and strict margin are universal finite certificate
conditions.  Under \(P=NP\), one first chooses one finite
\(z_0\in\R\), and only then makes finitely many existential successor choices.
\item[] \textbf{Model:} nonnegative joint time--space exponent vectors for
uniform exact SAT machines on one fixed explicit encoding and endpoint model.
The theorem itself is an ordered affine envelope; a future application must
prove model transport separately.
\item[] \textbf{SAME\_REALIZATION:} each \(z_j=(t_j,s_j)\) must be witnessed
jointly by one uniform exact SAT machine.  The phase \(q_j\) is theorem
bookkeeping, not free machine state or a license to mix coordinates.
\item[] \textbf{Weakest useful theorem:} finite graph-local potential
inequalities plus the displayed SAT-specific successor premise.  Neither a
common potential nor a universal source-code compiler is required.
\item[] \textbf{Falsifiers:} a negative matrix, missing threshold-to-forbidden
set implication, non-strict limiting margin, phase-erased illegal word,
coordinate mixing, uncharged successor construction, or any graph/depth
depending on SAT input length without a uniform algorithm proof.
\end{description}

\section{CORPUS REVIEWED}
The run surveyed the accessible repository tree and read root
\texttt{main.tex}, \texttt{main2.tex}, \texttt{main3.tex},
\texttt{main4.tex}, \texttt{main5.tex}, \texttt{AGENTS.md},
\texttt{CODEX.md}, the complete pnp-research skill and its mode/protocol
references, the proof-spines protocol, the current task index, the newest
local MAIN snapshot, governance material, all locally available LEAD
ARC reports, the current PROOF WRP-0005 and RESEARCH RSR-0005 reports, the
return-map specifications and finite-control programs, current root
additions, bibliographies, and cited raw experiments.  The preflight
completed with status OK.

The immediately preceding RESEARCH report was fetched from GitHub, was
substantive, and retained its unique beginning and final verification
sentinels at blob
\texttt{8d69f15446059b59313872057cc2ca4b79119c58f}.  Its current root writes
were independently fetched at blobs
\texttt{27daa948541caa62ba781c60aba97c58f5f2baa3} for
\texttt{main4.tex} and
\texttt{b1741807ce6d0bd81b86946e1dba1521e0fbc95e} for
\texttt{main5.tex}; no predecessor recovery was needed.

Repository searches for multinorm, path-complete, phase-dependent Lyapunov,
and copositive Lyapunov returned no prior project artifact.  Connected Drive
searches for the same terms returned no matching construction.  A broader
Drive search for phase dependent potential ARC recovered the older
Three-Regime ARC report; its bounded-versus-growing-depth distinction agrees
with the present quantifiers but contains no finite-state phase potential.

\section{DUPLICATION / CONTRADICTION AUDIT}
ARC-RETURN-CYCLE1 already proves finite entry for one fixed stable composite
\((P,q)\), including a two-matrix example in which neither individual phase
is stable.  COFINAL-SEED-HITTING1 further shows that a single finite route
from a suitable fixed seed is enough.  Those results are adopted rather than
relabelled.

PHASE-MULTINORM-ARC1 differs in its certificate and closure premise.  It
allows the selected edge and next phase to depend on the current achieved
pair, so there need not be one prescribed return word.  A positive linear
weight is attached to each graph node, and local edge inequalities telescope
without multiplying every possible word.  A deterministic cycle is a
special case, so the result extends rather than contradicts the old return
theorem.  It does not turn the old arithmetic example into a new SAT
transformation.

No verified root claim was strengthened.  The current density-gated
pigeonhole member is a partial VALUE cell, not a graph edge on every CNF.
STREAMCOMP1 and ORBITBATCH1 remain conditional access/accounting interfaces.
The new theorem supplies neither interface's missing all-input realization.

\section{Primary literature and exact transfer boundary}
Philippe, Essick, Dullerud, and Jungers study discrete-time switching systems
whose legal switching sequences are accepted by a fixed automaton
\cite{arc0007PhilippeMultinorm}.  Their Definition~2.2 assigns one norm per
automaton node and imposes an inequality for every labelled edge.
Proposition~2.3 identifies the infimum multinorm value with the constrained
joint spectral radius; Theorem~2.4 states that stability is equivalent to a
multinorm of value below one.  The model is finite-dimensional real linear
dynamics, the quantifier is over every accepted switching sequence, and the
resource is exponential state growth.

Classification: \textbf{THEOREM TRANSFER} for the finite graph, one-potential-
per-node, edge-inequality mechanism.  The ARC theorem below reproves the
needed positive affine form directly, including drift and forbidden-set
threshold, so no unproved control result is imported.  It is
\textbf{ANALOGY / CANDIDATE} for SAT closure: exponent pairs are not physical
states, achievable successors are existential uniform algorithms, and the
paper supplies no transformation from one SAT machine to another.

Ahmadi, Jungers, Parrilo, and Roozbehani introduce path-complete graph
Lyapunov functions for arbitrary switching and approximate the joint spectral
radius with graph-indexed inequalities
\cite{arc0007AhmadiPathComplete}.  This supports the multiple-potential
design space, but its path-completeness quantifies over all switching words.
ARC needs only one permitted existential successor at each finite proof stage.
Importing arbitrary-switching coverage or its semidefinite hierarchy would
therefore be a \textbf{FAILED TRANSFER} unless a SAT theorem supplies the
corresponding word set.

The post-quantum prompts do not match this obstacle.  Schur squares,
quadratic-relation spaces, syzygies, collision entropy, and binding extraction
do not produce an achievable successor edge.  Hash or random-oracle methods
would at best choose or index phases under an inadmissible assumption.
Classification for this run: \textbf{FAILED TRANSFER}.  A hidden finite
phase structure would in any event be algorithmically exploitable rather than
an automatic hardness source.

\section{Finite-state phase-potential theorem}
\begin{definition}[Phase-labelled positive affine closure]
Let \(G=(Q,E)\) be a fixed finite directed graph.  Each edge
\(e:q\to r\) has an entrywise nonnegative matrix
\(M_e\in\mathbb R_{\ge0}^{2\times2}\) and drift
\(d_e\in\mathbb R_{\ge0}^{2}\).  The achievable region has
\(G\)-closure if, for every \(q\in Q\) and
\(z\in\R\setminus\B\), some outgoing \(e:q\to r\) and some
\(z'\in\R\) satisfy
\[
 z'\le M_ez+d_e
\]
componentwise.
\end{definition}

\begin{theorem}[PHASE-MULTINORM-ARC1]
Assume \(\R\) contains a finite point and has \(G\)-closure.  Suppose there
are row vectors \(v_q\in\mathbb R_{>0}^2\), constants
\(0<\lambda<1\), \(c\ge0\), and \(\theta>0\) such that for every
\(e:q\to r\),
\[
 v_r^{\mathsf T}M_e\le\lambda v_q^{\mathsf T},
 \qquad v_r^{\mathsf T}d_e\le c,
 \tag{1}
\]
componentwise, and for every \(q\),
\[
 \{z\ge0:v_q^{\mathsf T}z<\theta\}\subseteq\B,
 \qquad \frac{c}{1-\lambda}<\theta.
 \tag{2}
\]
Then \(\R\cap\B\ne\varnothing\).
\end{theorem}
\begin{proof}
Choose any phase \(q_0\) and a finite \(z_0\in\R\).  While
\(z_j\notin\B\), \(G\)-closure supplies an edge
\(e_j:q_j\to q_{j+1}\) and \(z_{j+1}\in\R\) with
\(z_{j+1}\le M_{e_j}z_j+d_{e_j}\).  Put
\(V_j=v_{q_j}^{\mathsf T}z_j\).  Positivity of \(v_{q_{j+1}}\) and (1) give
\[
\begin{aligned}
V_{j+1}
&\le v_{q_{j+1}}^{\mathsf T}(M_{e_j}z_j+d_{e_j})\\
&\le\lambda v_{q_j}^{\mathsf T}z_j+c
 =\lambda V_j+c.
\end{aligned}
\]
Therefore
\[
 V_j\le\lambda^jV_0+c\frac{1-\lambda^j}{1-\lambda}.
\]
The right side converges to \(c/(1-\lambda)<\theta\).  Some finite \(j\)
has \(V_j<\theta\), and (2) gives \(z_j\in\B\).
\end{proof}

The theorem has a finite certificate once the SAT-successor premise is
available: all matrix and drift inequalities are linear.  It does not assert
that linear phase potentials are necessary, nor that a stable SAT relation
must expose rational weights of bounded description.

For the Williams half-plane, a convenient threshold follows from
\[
 m=\min_{q\in Q}\min\{(v_q)_1,(v_q)_2\}>0.
\]
Since \(v_q^{\mathsf T}(t,s)\ge m(t+s)\), any
\(\theta\le m\gamma\) satisfies the first part of (2).  This is the exact
threshold-to-forbidden-set bridge missing from a bare claim that a potential
decreases.

\section{Strict improvement over one common linear potential}
Consider phases \(A,B\), weights and constants
\[
v_A=(1,1),\qquad v_B=(1/3,3),\qquad
\lambda=2/3,\quad c=1/100,\quad\theta=1/3.
\]
The two cross edges use
\[
M_{AB}=\begin{pmatrix}2&0\\0&1/5\end{pmatrix},\quad
d_{AB}=\binom{1/100}{0},
\]
\[
M_{BA}=\begin{pmatrix}1/5&0\\0&2\end{pmatrix},\quad
d_{BA}=\binom{0}{1/300}.
\]
Both phases also have self-loops with matrix \((1/2)I\) and drift
\((1/600,1/600)^{\mathsf T}\).  Exact multiplication gives
\[
\begin{array}{c|c|c}
e & v_r^{\mathsf T}M_e & v_r^{\mathsf T}d_e\\
\hline
A\to A&(1/2,1/2)&1/300\\
A\to B&(2/3,3/5)&1/300\\
B\to A&(1/5,2)&1/300\\
B\to B&(1/6,3/2)&1/180
\end{array}
\]
and every row is bounded by
\(\lambda v_q^{\mathsf T}\) with drift at most \(c\).
Moreover \(c/(1-\lambda)=3/100<1/3\).
Because the smallest weight coordinate is \(1/3\),
\(v_q^{\mathsf T}(t,s)<1/3\) implies \(t+s<1\); and
\(\pi/7<\pi/3\) gives \(1<2\cos(\pi/7)\).

No common positive row vector \(w\) can contract every edge matrix with any
\(\mu<1\): applying
\(w^{\mathsf T}M_{AB}\le\mu w^{\mathsf T}\) to the first coordinate would
give \(2w_1\le\mu w_1\), impossible for \(w_1>0\).
Thus phase labels strictly improve the class of \emph{linear-potential
certificates}.  This does not prove a larger class of SAT successors exists;
the entire model is arithmetic.

\section{Reproducible exact-rational test}
The command
\begin{verbatim}
timeout 30s python3 experiments/proof_spines/phase_multinorm_certificate.py \
  results/raw/proof_spines/2026-10-08_063447_phase_multinorm.json
\end{verbatim}
uses a 25-second cooperative alarm, seed 20261008, and
\texttt{fractions.Fraction} throughout.  It completed in 22.288 seconds.
For each initial phase it enumerated all \(2^{10}=1024\) accepted graph
words of depth ten on 65 rational starting points, checking 1,331,200
individual transition and envelope inequalities.  It observed 132,196
threshold crossings within depth ten; this finite count is
\textsc{Computational} and is not used to prove asymptotic crossing.

The code SHA-256 is
\path{6063e00b010c9bc10994a0c04fee2313047a1d9917f75856a2077fb9a5cd4894};
the raw JSON SHA-256 is
\path{5bf6c971aee9f3b544b89f5a824ca67a3b08cde9bd5402934169f5db524e9f97}.

\begin{description}
\item[] \textbf{Common-potential control.} The unstable \(2\) diagonal entry
proves exact impossibility of a common positive linear contraction.
\item[] \textbf{Phase-erasure control.} Repeating \(M_{AB}\) would grow the
first coordinate as \(2^k\), but the second copy is not a legal edge out of
phase \(B\).  The graph label is load-bearing.
\item[] \textbf{Negative-matrix control.} With
\(M=\operatorname{diag}(-1,1/2)\), \(z=(0,0)\le u=(1,0)\), but
\(Mz=(0,0)\not\le(-1,0)=Mu\).  Nonnegativity is load-bearing.
\item[] \textbf{Easy-family and independent-product control.} No CNF,
restriction, or SAT solver is instantiated.  Horn, 2-SAT, affine/Tseitin,
bounded-treewidth, and product families cannot promote this model into an
all-CNF successor theorem.
\item[] \textbf{Source/current-scale control.} No source-carry or SAT call is
hidden in the arithmetic.  A real successor must charge both if present.
\item[] \textbf{Padding control.} Every coordinate is an exponent on the
original fixed encoding.  Rebaselining against padded length is not an edge.
\item[] \textbf{Oracle/algebraic control.} The envelope proof survives
semantic or algebraic oracle replacement whenever its premise is assumed.
It is therefore not the nonrelativizing or nonalgebrizing hinge.
\end{description}

\section{Finite existence versus growing-depth uniformity}
The phase graph does not create one algorithm that traverses infinitely many
states.  Under \(P=NP\), first fix one finite achievable exponent pair.
The theorem proof then chooses a finite list
\[
(q_0,z_0),(q_1,z_1),\ldots,(q_J,z_J),
\]
where every \(z_j\) is witnessed by one uniform exact SAT machine and
\(J\) is a constant after \(z_0\) has been fixed.  Machines and edges may
differ across the finite proof.  No procedure must accept arbitrary source
code and emit the next machine.

By contrast, allowing \(Q=Q_N\), \(\lambda=\lambda_N\to1\), or
\(J=J(N)\to\infty\) as SAT input length grows defines a single
growing-depth uniform algorithm problem.  It requires an effective scheduler,
uniform \(o(1)\) rates, description access, construction cost, and live-space
accounting.  PHASE-MULTINORM-ARC1 proves none of these.  A finite automaton
whose size is merely finite for each \(N\) is not a fixed finite certificate.

\section{Barrier audit and remaining implication}
The theorem is order-theoretic and relativizing.  DEAD, COPY, DELAY, INTERP,
RECOMP, semantic-oracle, and algebraic-oracle implementations can bypass or
inflate any proposed SAT edge unless the successor theorem charges them on
the same realization.  Natural Proofs is not implicated because no
constructive property of Boolean functions is used.

The load-bearing statement remains
\[
\boxed{\begin{gathered}
P=NP\Longrightarrow \exists\text{ a fixed finite graph and certified
phase potentials}\\
\text{such that every outside achievable SAT pair has one graph-legal,
jointly witnessed successor.}
\end{gathered}}
\]
Its matrix inequalities are now a solved conditional interface.  The
SAT-specific existence implication is \textsc{Open}.  One open arrow means
there is no separation proof.

\section{WHAT CHANGED RELATIVE TO ROOT MAIN}
Root \texttt{main.tex} was reviewed and left unchanged because this is not an
LSPV ownership-region run and the user assigns ARC/WRAPPER detail to
\texttt{main4.tex} and \texttt{main5.tex}.  The ARC spine gains the
phase-dependent theorem, exact rational certificate, transfer boundary, and
finite-versus-growing-depth guard.  The WRAPPER spine receives a concise
review stating that no existing component realizes the new premise.

\section{ADOPTED / REJECTED / RETIRED PRIOR COMPONENTS}
\begin{description}
\item[] \textbf{Adopted:} ARC-RETURN-CYCLE1, cofinal finite seed hitting,
entrywise-nonnegative order propagation, same-realization exponent pairs, and
the fixed finite versus input-growing-depth distinction.
\item[] \textbf{Adopted as exact envelope machinery:} finite-state
multinorm edge inequalities, with drift and an explicit forbidden-set
threshold added here.
\item[] \textbf{Rejected:} erasing phase labels; importing arbitrary-switching
quantifiers; treating a control state as a free SAT algorithm; or claiming
that potential feasibility constructs a SAT successor.
\item[] \textbf{Retired:} no prior proved component is retired.  A common
positive linear potential is only a sufficient special case and is inadequate
for the displayed arithmetic system.
\end{description}

\section{STALE OR OVERSTATED CLAIMS FOUND}
The archived ARC text says common Lyapunov functions are only a control-theory
analogy.  That sentence is not false, but it is incomplete: the finite
graph-local potential calculation transfers exactly after the SAT-successor
relation is separately assumed.  The new root addition records this narrower
exact transfer without promoting the missing SAT bridge.  No existing
\textsc{Open} label was stale.

\section{Peer handoff}
\textbf{PROOF/RED target:} take one concrete original-SAT-scale portfolio
with two genuinely different joint time--space ledgers and state its legal
phase graph.  Either produce exact nonnegative edge bounds and positive
weights satisfying (1)--(2), with failed acquisition and source access
charged, or minimize a formula/implementation on which every outgoing edge
fails.  Do not search for another abstract potential until one general-CNF
successor edge is source-visible.

\bibliographystyle{plain}
\bibliography{spines,arc_orbit_refs}
\end{document}
% END_OF_VERIFIED_REPORT::LEAD::ARC-0007
